% Copyright
      \vspace{0.5cm}

UNIVERSIDADE FEDERAL DO RIO DE JANEIRO \\
Escola Politécnica - Departamento de Engenharia Mecânica \\
Centro de Tecnologia, bloco G, sala G-204, Cidade Universitária \\ 
Rio de Janeiro - RJ      CEP 21941-909\\
\vspace{0.5cm}
\paragraph{}Este exemplar é de propriedade da Universidade Federal do Rio de Janeiro, que poderá incluí-lo em base de dados, armazenar em computador, microfilmar ou adotar qualquer forma de arquivamento.
\paragraph{}É permitida a menção, reprodução parcial ou integral e a transmissão entre bibliotecas deste trabalho, sem modificação de seu texto, em qualquer meio que esteja ou venha a ser fixado, para pesquisa acadêmica, comentários e citações, desde que sem finalidade comercial e que seja feita a referência bibliográfica completa.
\paragraph{}Os conceitos expressos neste trabalho são de responsabilidade do(s) autor(es).


\pagebreak

% Dedicatória
\begin{center}
\textbf{DEDICATÓRIA}
\end{center}
      \vspace{0.5cm}

\paragraph{}Dedico este trabalho, primeiramente, a Deus, Senhor da minha vida, que sustentou cada etapa desta caminhada e me concedeu graça, perseverança e propósito quando as forças pareciam insuficientes.

\paragraph{}À minha esposa Aretha, companheira fiel, amiga e apoio constante em todos os momentos, cujo amor, paciência e incentivo tornaram possível a conclusão desta jornada.

\paragraph{}À minha filha Marina, que mesmo sem compreender totalmente o significado deste trabalho, foi diariamente a maior motivação para não desistir. Que este diploma seja também um testemunho para você de que vale a pena terminar o que começamos.

\paragraph{}Aos meus pais, Daniela e Silvio, que sempre acreditaram na educação como instrumento de transformação e nunca deixaram faltar apoio, cuidado e exemplo.

\paragraph{}Às minhas irmãs, Yngrid e Yasmim, pela presença, amizade e pelo vínculo familiar que sempre me fortaleceu.

\vspace{1.5cm}
\begin{center}
\textit{Comecei este caminho como estudante. Concluo como engenheiro, marido, pai e servo de Cristo.}
\end{center}

\pagebreak


% Agradecimento
\begin{center}
\textbf{AGRADECIMENTO}
\end{center}
      \vspace{0.5cm}

\paragraph{}Agradeço primeiramente a Deus, por sustentar cada etapa desta caminhada e por conceder sabedoria e perseverança para chegar até aqui.

\paragraph{}Ao meu orientador, Prof. Gustavo Rabello dos Anjos, pela orientação cuidadosa, confiança e incentivo intelectual, fundamentais para o desenvolvimento deste trabalho e para o amadurecimento da proposta do MinervaMesh.

\paragraph{}À Universidade Federal do Rio de Janeiro e ao corpo docente do curso de Engenharia Mecânica, pela formação técnica e científica que tornou possível a realização deste projeto.

\paragraph{}Ao meu grande amigo e padrinho de casamento, Michel Silva Bonifácio, pelo apoio durante a graduação e pela presença constante em momentos decisivos da minha trajetória acadêmica.

\paragraph{}Ao meu amigo Eliab Araújo, pelo incentivo para a conclusão do curso e pela motivação contínua no desenvolvimento deste trabalho.

\paragraph{}Aos colegas, amigos e a todos que contribuíram direta ou indiretamente para minha formação como engenheiro.

\pagebreak


% Resumo
\begin{center}
\textbf{RESUMO}
\end{center}
      \vspace{0.5cm}

\paragraph{}A simulação numérica é ferramenta central de análise e projeto em Engenharia Mecânica, mas seu uso no ensino de métodos numéricos esbarra em duas dificuldades recorrentes: as ferramentas mais capazes exigem licenciamento, infraestrutura dedicada e configuração de ambiente, e raramente expõem de forma transparente a formulação matemática e o código por trás de cada simulação. O problema atacado neste trabalho é, portanto, reduzir a barreira de acesso a um ambiente de simulação numérica didático, sem abrir mão do rigor técnico nem da transparência da implementação. Como resposta, apresenta-se o \textit{MinervaMesh}, uma plataforma que resolve Equações Diferenciais Parciais de interesse em Engenharia Mecânica pelo Método dos Elementos Finitos (MEF) e pelo Método das Diferenças Finitas (MDF). A contribuição articula dois aspectos complementares, subordinados a esse único objetivo: (i) a implementação e a verificação dos métodos numéricos, confrontadas prioritariamente com soluções analíticas conhecidas e, nos casos sem solução fechada, com \textit{benchmarks} consagrados da literatura; e (ii) a entrega dessas simulações integralmente no navegador, em arquitetura \textit{client-side} sobre \textit{WebAssembly} e \textit{PyScript}, sem instalação e com o código-fonte Python inspecionável em cada caso. Os resultados mostram que as equações discretizadas reproduzem corretamente as soluções de referência e que a plataforma é adequada ao ensino, à prototipagem e a demonstrações acadêmicas de mecânica computacional.
\paragraph{}
\noindent Palavras-Chave: Método dos Elementos Finitos, Método das Diferenças Finitas, Simulação Web, Mecânica Computacional, Ensino de Engenharia.

\pagebreak


% Abstract
\begin{center}
\textbf{ABSTRACT}
\end{center}
      \vspace{0.5cm}

\paragraph{}Numerical simulation is a core tool for analysis and design in Mechanical Engineering, yet teaching numerical methods runs into two recurring obstacles: the most capable tools require licensing, dedicated infrastructure, and environment setup, and they seldom expose the underlying mathematical formulation and source code of each simulation. The problem addressed in this work is therefore to lower the barrier of access to a didactic numerical simulation environment, without sacrificing technical rigor or implementation transparency. As a response, this work presents \textit{MinervaMesh}, a platform that solves Partial Differential Equations of interest in Mechanical Engineering with the Finite Element Method (FEM) and the Finite Difference Method (FDM). The contribution combines two complementary aspects, subordinated to this single goal: (i) the implementation and verification of the numerical methods, assessed primarily against known analytical solutions and, for cases without a closed-form solution, against well-established benchmarks; and (ii) the delivery of these simulations fully in the browser, using a \textit{client-side} architecture built on \textit{WebAssembly} and \textit{PyScript}, with no installation and with inspectable Python source code for every case. Results show that the discretized equations correctly reproduce the reference solutions and that the platform is suitable for teaching, rapid prototyping, and academic demonstrations in computational mechanics.
\paragraph{}
\noindent Key-words: Finite Element Method, Finite Difference Method, Web Simulation, Computational Mechanics, Engineering Education.

\pagebreak


% Siglas
\begin{center}
\textbf{SIGLAS}
\end{center}
      \vspace{0.5cm}

\paragraph{}UFRJ - Universidade Federal do Rio de Janeiro
\paragraph{}MEF - Método dos Elementos Finitos
\paragraph{}MDF - Método das Diferenças Finitas
\paragraph{}CFD - \textit{Computational Fluid Dynamics} (Dinâmica dos Fluidos Computacional)
\paragraph{}WASM - WebAssembly
\paragraph{}SUPG - \textit{Streamline Upwind/Petrov-Galerkin}


\pagebreak






